% sections/12_extended_hypotheses.tex — H3, H4, limitations (v3)

\section{Extended Analysis: New Observations}

The dataset's severity definition --- measuring traffic disruption rather than
physical harm --- opens several additional hypotheses not captured by the
original two.

%% ── H3 ──────────────────────────────────────────────────────
\subsection{O1: Infrastructure Controls Reduce Severity More Than Weather Does}

\begin{quote}
  \itshape
  Road infrastructure with mandatory stopping controls (signals, crossings,
  stop signs) reduces accident severity more substantially than any single
  weather category increases it.
\end{quote}

\subsubsection*{Evidence from V7}

Traffic signals reduce the high-severity rate from 21.19\% to 9.51\% --- a
difference of \textbf{11.7\,pp}. By comparison, the worst weather category
(Rain) raises the severity rate by only 6.53\,pp above Clear (24.38\%
vs.\ 17.85\%) after outlier removal.

\begin{equation*}
  \Delta_{\text{infrastructure (signal)}} = 11.7\,\text{pp}
  \qquad \gg \qquad
  \Delta_{\text{weather (Rain vs Clear)}} = 6.53\,\text{pp}
\end{equation*}

Infrastructure intervention has roughly $\mathbf{1.8\times}$ the effect
magnitude of the worst adverse weather on severity.

\subsubsection*{Why This Matters}

Since severity tracks traffic disruption duration, infrastructure that organises
traffic (signals, crossings) ensures faster clearance, fewer secondary
collisions, and quicker emergency access --- all reducing the severity score
directly.

\verdict{O1}{Supported}{%
  The infrastructure effect size dominates the weather effect size even after
  outlier removal tightened the weather signal. Signal installation ($-11.7$\,pp)
  outperforms the strongest weather comparison ($+6.53$\,pp for Rain vs.\ Clear).
  Crossing removal ($-14.0$\,pp) and stop sign removal ($-13.6$\,pp) show even
  larger gaps.%
}

%% ── H4 ──────────────────────────────────────────────────────
\subsection{O2: High-Speed, Low-Traffic Conditions Produce Disproportionately
High Severity}

\begin{quote}
  \itshape
  Accidents during low-traffic periods (night, early morning, weekends) produce
  higher per-accident severity scores due to higher vehicle speeds and longer
  emergency response times.
\end{quote}

\subsubsection*{Evidence from V6a and V6b}

After outlier removal, the early morning plateau (hours 3--6) is cleaner and
more clearly sustained:

\begin{table}[H]
\centering\small
\begin{tabular}{@{}cc@{}}
  \toprule
  \hdrrow \textbf{Hour} & \textbf{High-Severity Rate} \\ \midrule
  3\,h & \todobox{add value} \\
  4\,h & 21.05\% \\
  5\,h & 21.49\% \\
  6\,h & 21.73\% \\
  \midrule
  \midrow \textit{Rush 7--8\,am} & \textit{18.83--18.99\% (below average)} \\
  \bottomrule
\end{tabular}
\caption{Early morning severity plateau (hours 3--6) vs.\ rush hours}
\end{table}

The Saturday 7--9\,am darkest cells on V6b (peak $\approx$0.3176) remain the
single strongest temporal interaction, confirming the weekend-low-traffic effect.

\subsubsection*{Evidence from V5d (Visibility Analysis)}

Very High visibility ($>10$\,mi) accidents show a 31.65\% high-severity
rate --- higher than any other visibility bin. High visibility correlates with
rural highways and open roads where speeds are elevated, reinforcing the
speed-severity link.

\subsubsection*{Corroboration from V\_B2}

Oregon's wide IQR in V\_B2 reflects genuine longer clearance times on rural
highways where emergency response is slower --- exactly the mechanism proposed
by H4. The severity score in this dataset is fundamentally a \emph{clearance
time} score, and clearance time is longest when speed is high and responders
are far away.

\verdict{O2}{Supported}{%
  Low-traffic, high-speed conditions are the primary driver of high per-accident
  severity. The signal is stronger and cleaner after $2\times$IQR removal: the
  early morning plateau (hours 3--6) is sustained, and the Saturday 7--9\,am
  heatmap peak is the single strongest interaction across all temporal views.%
}

%% ── Limitations ─────────────────────────────────────────────
\subsection{Limitations and Paths for Deeper Analysis}

\begin{itemize}[noitemsep]
  \item \textbf{Speed Data Integration:} The dataset lacks posted speed limits.
        Merging with NHTSA speed data by road segment would directly test H4.

  \item \textbf{Multivariate Severity Modelling:} A logistic regression or
        decision tree on \texttt{Weather\_Group}, \texttt{Road\_Feature},
        \texttt{Hour}, and \texttt{Is\_Weekend} simultaneously would reveal
        interaction effects invisible in single-variable charts.

  \item \textbf{Duration as Dependent Variable:} Using \texttt{Duration\_Min}
        rather than \texttt{Severity\_Group} as the outcome leverages the
        continuous signal for finer-grained analysis.

  \item \textbf{State-Level Heterogeneity:} California contributes 23\% of
        records. State-stratified analysis would reveal whether H1 and H2 hold
        uniformly or are driven by California's climate and road network.

  \item \textbf{COVID Counterfactual:} 2020--2021 saw 30--60\% lower traffic
        volume with maintained infrastructure. Comparing pre/post COVID severity
        rates isolates volume effects from infrastructure effects.

\end{itemize}
